% Sumber: disusun dari Bab 4 tesis ini (EXP-005 sampai EXP-011); tidak ada klaim baru yang diperkenalkan di sini.
\chapter{Kesimpulan}
\label{chap:conclusion}

\section{Jawaban atas Pertanyaan Penelitian}

\textbf{RQ1.} Pada subset 135 potongan, akurasi HSEmotion adalah 0,483 pada kelompok Gelap dan 0,130 pada kelompok Gelap-Sedang, dengan kesenjangan 0,353 (SK 95\% [0,182; 0,540]). Kesenjangan ini tetap terlihat di dalam kelas Netral dan Senang, dan sebagian besar berasal dari potongan Netral Gelap-Sedang yang diprediksi Marah (12 dari 21). Arah yang sama teramati pada sampel 227 potongan. Ini adalah asosiasi pada satu sampel, bukan bukti bahwa warna kulit menjadi penyebabnya.

\textbf{RQ2.} Konfigurasi CLAHE adaptif yang diuji mengubah kesenjangan dari 0,353 menjadi 0,330 (perubahan $-$0,023; SK 95\% [$-$0,136; 0,073]). Pengaruhnya terhadap kesenjangan tidak dapat disimpulkan dari data ini; hasil ini tidak membuktikan bahwa CLAHE tidak berpengaruh. Kesenjangan tetap besar dengan maupun tanpa CLAHE, dan laju prediksi Marah pada Gelap-Sedang tidak berubah (0,478).

\textbf{RQ3.} ArcFace + Regresi Logistik memiliki kesenjangan yang lebih kecil (0,123; SK 95\% [$-$0,038; 0,282]) karena akurasinya pada Gelap-Sedang lebih tinggi (0,326 berbanding 0,130), sementara akurasi Gelap-nya tidak lebih baik. Selisih kesenjangan kedua model tidak memuat nol pada pengelompokan 135 potongan (0,229; SK 95\% [0,054; 0,442]) tetapi memuat nol pada pengelompokan 78 potongan (0,188; SK 95\% [$-$0,083; 0,477]), sehingga pengurangan kesenjangan ini bersifat indikatif, belum terbukti. Akurasi keseluruhan kedua model tidak berbeda secara signifikan (McNemar $p = 0,4614$), dan perbandingan ini bersifat retrospektif.

\section{Ringkasan}
\label{sec:summary}

Tesis ini mengevaluasi kesenjangan kinerja dua pendekatan pengenalan ekspresi wajah antara dua kelompok kecerahan kulit pada potongan wajah dari satu film dokumenter. Kesenjangan HSEmotion terukur dengan jelas pada sampel ini; satu konfigurasi praproses kontras tidak menghasilkan perubahan yang dapat disimpulkan; dan satu klasifier alternatif menunjukkan kesenjangan yang lebih kecil dengan bukti yang masih indikatif. Kedua kelompok sama-sama gelap, sehingga tesis ini tidak menjawab pertanyaan tentang kinerja pada kulit terang dibanding kulit gelap. Hasil ini tidak berlaku umum di luar sampel yang dievaluasi.

\section{Kontribusi}

Kontribusi tesis ini adalah: (1) evaluasi kesenjangan antar-kelompok kecerahan yang dapat direproduksi dari artefak tersimpan (EXP-010 dan EXP-011), lengkap dengan selang kepercayaan \textit{bootstrap} berkelompok; (2) dokumentasi bahwa pola kesalahan Netral$\to$Marah HSEmotion terkonsentrasi pada kelompok Gelap-Sedang dan menjelaskan sebagian besar kesenjangannya pada sampel ini; dan (3) pelaporan yang hati-hati atas dua intervensi sederhana, praproses dan penggantian klasifier, yang membedakan hasil yang tidak dapat disimpulkan dan hasil yang masih indikatif dari hasil yang sudah terbukti.

\section{Keterbatasan}

Keterbatasan dirinci pada Bab~\ref{chap:results} \S\ref{sec:results-limitations}. Yang paling membatasi klaim adalah: kelompok kecerahan sebagai proksi yang bercampur dengan pencahayaan dan hanya mencakup dua kelompok gelap; kelompok Gelap-Sedang yang kecil (46 potongan); satu film sumber dan satu anotator tanpa ukuran kesepakatan; ketiadaan identitas subjek; sifat retrospektif RQ3; dan dasar persetujuan serta asal-usul rekaman sumber yang belum terdokumentasi.

\section{Kerja Lanjutan}
\label{sec:future-work}

\begin{itemize}
  \item \textbf{Representasi kelompok.} Menambah potongan Gelap-Sedang serta kelompok yang lebih terang, dengan protokol perekaman dan persetujuan yang jelas, agar kesenjangan dapat diukur lebih presisi dan perbandingan kulit terang dan gelap menjadi mungkin.
  \item \textbf{Kualitas label.} Melibatkan anotator kedua dan melaporkan kesepakatan antar-anotator, termasuk menyelesaikan 75 potongan tumpang-tindih Papuan V2.
  \item \textbf{Validitas proksi.} Memeriksa secara manual sebagian kelompok kecerahan untuk menilai seberapa jauh kelompok tersebut dipengaruhi pencahayaan.
  \item \textbf{Pembagian data sadar-identitas.} Membangun identitas subjek atau jejak wajah agar partisi latih dan uji tidak berbagi orang yang sama.
  \item \textbf{Uji yang dirancang di muka.} Menguji pendekatan untuk memperkecil kesenjangan, misalnya klasifier berbobot kelompok, dengan kriteria keberhasilan yang ditetapkan sebelum evaluasi.
  \item \textbf{Privasi dan etika.} Menyelesaikan dasar persetujuan penggunaan rekaman, dan hanya memublikasikan hasil agregat serta ilustrasi wajah yang sudah dianonimkan.
  \item \textbf{Titik-fitur wajah.} Jalur titik-fitur tetap pekerjaan pendukung: 64 gambar sudah diperiksa manusia, tetapi belum ada model titik-fitur yang terbukti lebih baik daripada \textit{mean-shape} (\S\ref{sec:methodology-landmarks}), dan belum ada eksperimen terkontrol yang menghubungkan titik-fitur dengan kinerja klasifikasi emosi atau kesenjangannya.
\end{itemize}
